\documentclass[11pt,a4paper]{article}
\usepackage[T1]{fontenc}
\usepackage[utf8]{inputenc}
\usepackage{lmodern,microtype}
\usepackage[margin=27mm,headheight=15pt]{geometry}
\usepackage{amsmath,amssymb,amsthm,mathtools}
\usepackage{booktabs,array,longtable,enumitem}
\usepackage[dvipsnames]{xcolor}
\usepackage{graphicx,fancyhdr}
\usepackage{hyperref}
\hypersetup{colorlinks=true,linkcolor=MidnightBlue,citecolor=MidnightBlue,
 urlcolor=MidnightBlue,pdftitle={The source boundary of extensional acyclic digraphs},
 pdfauthor={Research prepared for Vladimir Reshetnikov},
 pdfsubject={OEIS A182220, A182162, A001192; extremal enumeration and non-holonomicity}}
\pagestyle{fancy}\fancyhf{}
\fancyhead[L]{\small\textsc{The source boundary of extensional DAGs}}
\fancyhead[R]{\small OEIS research}
\fancyfoot[C]{\thepage}
\renewcommand{\headrulewidth}{0.3pt}
\newtheorem{theorem}{Theorem}[section]
\newtheorem{proposition}[theorem]{Proposition}
\newtheorem{lemma}[theorem]{Lemma}
\newtheorem{corollary}[theorem]{Corollary}
\theoremstyle{definition}
\newtheorem{definition}[theorem]{Definition}
\newtheorem{example}[theorem]{Example}
\newtheorem{remark}[theorem]{Remark}
\newcommand{\N}{\mathbb N}
\newcommand{\PP}{\mathcal P}
\newcommand{\U}{\mathcal U}
\newcommand{\E}{\mathbb E}
\newcommand{\Var}{\operatorname{Var}}
\newcommand{\TV}{d_{\mathrm{TV}}}
\newcommand{\sour}{\operatorname{src}}
\newcommand{\bits}{\lceil\log_2 n\rceil}
\newcommand{\fall}[2]{(#1)_{\underline{#2}}}
\newcommand{\B}{\mathcal B}
\newcommand{\code}[1]{\texttt{\detokenize{#1}}}
\DeclareMathOperator{\Bin}{Bin}
% Added by the ProveIt write phase (batch 85): repository paths in the
% [write] notes may also break after a hyphen; dated [write] notes.
\expandafter\def\expandafter\UrlBreaks\expandafter{\UrlBreaks\do\-}
% [write] typewriter text prints "--" as two hyphens (the delivered PDF set
% the options --max-n and --brute-n with an en dash).
\DisableLigatures[-]{encoding = T1, family = tt*}
\newenvironment{writenote}[1][]{\par\smallskip\begingroup\small
 \noindent\textbf{[write]}\if\relax\detokenize{#1}\relax\else~\textbf{#1.}\fi\ }%
 {\par\endgroup\smallskip}
\allowdisplaybreaks[2]
\setlist{nosep,leftmargin=1.6em}
\setlength{\parskip}{0.25em}

\begin{document}
\hypersetup{pageanchor=false}
\begin{titlepage}
\thispagestyle{empty}
\vspace*{8mm}
{\color{MidnightBlue}\large\textsc{Research article \quad / \quad October 3, 2026}}\par
\vspace{13mm}
{\LARGE\bfseries The source boundary of\par extensional acyclic digraphs}\par
\vspace{6mm}
{\Large An OEIS conjecture, finite-defect enumeration,\par and dyadic non-holonomicity}\par
\vspace{10mm}
{\large Prepared for Vladimir Reshetnikov}\par
\vspace{3mm}
{\normalsize In the mathematical research context of the ProveIt repository}
\vspace{12mm}
\begin{abstract}
An extensional acyclic digraph has pairwise distinct out-neighborhoods.
We give a self-contained proof of the formula
$a(n)=n-\lceil\log_2n\rceil$ for its maximum possible number of sources,
the formula still marked conjectural in OEIS A182220 when consulted.
The pigeonhole upper bound is already present in Tomescu's 2011 thesis;
we credit that fact and do not present it as a new discovery.

The main development concerns the boundary of the source-count triangle
A182162. If $u_m=\mathrm{A001192}(m)$, $q=\lceil\log_2 n\rceil$, and
$b_n^{(d)}$ counts isomorphism classes with exactly $n-q-d$ sources,
we obtain an exact $(d+1)$-term formula and, in particular,
$b_n^{(0)}=u_q\binom{2^q-q}{n-q}$.
A uniform covering estimate then proves a dyadic entropy law for every
fixed $d$. The complete set of subsequential exponential growth rates
is an explicit interval, equal to $[1,4]$ when $d=0$.
The exponential jumps across powers of two imply that none of these
boundary sequences is P-recursive. Consequently their ordinary generating
functions are not D-finite.

We also give an all-orders expansion of the Boolean-subset factor,
retaining the small-core count exactly, and a total-variation binomial
law for the number of missing arcs near a complete Boolean layer.
The proofs are accompanied by independently implemented counting
recurrences and exhaustive small-graph checks. These are conventional
mathematical proofs with computational tests, not Lean-verified results.
Historical priority for the boundary refinements is not established.
\end{abstract}
\vfill
\noindent\small
\textbf{Primary sequences:} A182220, A182162, A001192, A182161.\par
\textbf{Keywords:} transitive finite sets; source enumeration; Boolean layers;
finite defects; entropy; polynomial recurrences; exact sampling.
\end{titlepage}
\hypersetup{pageanchor=true}

{\small\setlength{\parskip}{0pt}\tableofcontents}
\newpage

\section{Problem, contribution, and prior-art boundary}

\subsection{What is being answered}
Let $a(n)$ be the largest possible number of sources in an extensional
acyclic digraph on $n\ge1$ vertices. OEIS A182220 identifies this quantity
and records, among other equivalent conjectural expressions,
\begin{equation}\label{sbd:eq:oeis-conjecture}
 a(n)=n-1-\lfloor\log_2(n-1)\rfloor\qquad(n\ge3).
\end{equation}
The entry attributes this formula to Antti Karttunen, August 17, 2013
\cite{oeis220}. For $n\ge2$, the right side equals
$n-\lceil\log_2n\rceil$; the latter also gives $a(1)=1$.
The problem is not an issue of numerical fitting. We prove both the
universal obstruction and an explicit construction for every admissible
source count.

The more substantial question is: \emph{how many graphs lie on, or a
fixed distance below, this extremal boundary, and how do those counts
behave?} Write $u(n,k)$ for the number of isomorphism classes with $n$
vertices and $k$ sources. The labeled count is $n!u(n,k)$; the factor
$n!$ will be justified, not assumed. Define
\begin{equation}\label{sbd:eq:boundary-def}
 q=q(n)=\lceil\log_2n\rceil,
 \qquad b_n^{(d)}=u(n,n-q-d),\qquad d\in\N_0,
\end{equation}
with value zero if the requested number of sources is not positive.
For each fixed $d$, this convention affects only finitely many terms.

\subsection{What was already known}
Peddicord studied the numbers $u_m$ of finite transitive, or full, sets
\cite{peddicord,oeis1192}. Policriti and Tomescu developed source-sensitive
enumeration of extensional acyclic digraphs \cite{policriti}.
Tomescu's thesis explicitly gives the source bound
$k\le n-\lceil\log_2n\rceil$ on printed page 29, along with source-counting
recurrences and rigidity results \cite{tomescu}. The source sieve used
below is also encoded in Johnston's 2012 Maple program in A182162
\cite{oeis162}. Wagner studied asymptotic enumeration and distributions
in the unrestricted graph class \cite{wagner}.

Accordingly, neither the upper bound, the correspondence with full sets,
rigidity, nor the general source sieve is claimed as new. An OEIS
conjectural label does not establish that a statement is an unsolved
research problem in the literature. Our proof supplies a direct resolution
of the displayed formula; its upper-bound half has explicit prior art.

The exact boundary classification, the finite-defect specialization,
the uniform entropy profile, and the non-P-recursiveness theorem form
the development of this article. Some exact identities are short
consequences of the classical source sieve. We make no assertion of
historical priority for those consequences or for the probabilistic
refinements. The meaningful claim is that complete proofs are supplied
here, not that an exhaustive literature search has certified novelty.

\subsection{Repository relationship and reproducibility}
The supplied ProveIt repository was inspected at commit
\begin{center}
\small\code{6bf7f30d0352f7596e70928b3d4f304914075907}.
\end{center}
Its organization separates mathematical sources, research reports, and
executable certificates \cite{repo}. Searches for the present OEIS
identifier did not locate an existing treatment. That limited search is
not an exhaustive duplication audit. No repository theorem is used as an
unverified premise below, and no repository files were modified.

The accompanying \code{verify.py} uses exact integers and independently
implements a marked-source sieve and a source-deletion recurrence. It
also enumerates all extensional DAG isomorphism classes through seven
vertices. The finite tests support error detection; the proofs, rather
than the tests, establish the infinite statements.

\begin{writenote}[Added 3 October 2026, batch~85]
\emph{Provenance.} This report is manuscript~07 of batch~85 of ProveIt's
\nolinkurl{docs/incoming} intake: the archive
\nolinkurl{OEIS_A182220_Source_Boundary_Research.zip} (wrapper directory
\nolinkurl{oeis_source_boundary/}), arrival commit \nolinkurl{9d6968c8a}, placed in
\nolinkurl{ddf8df5d5} (batch~85C), written against the pin \nolinkurl{6bf7f30d0}
quoted above. It is a single-source report. Its title block names no
author and no tool. The write phase prefixed every label with
\nolinkurl{sbd:}, added the dated notes marked [write], and kept the title page
out of the PDF page anchors; no statement, proof or number was changed.

\emph{Priority.} The formula $a(n)=n-\lceil\log_2n\rceil$ is \emph{not}
presented as a new discovery. Its upper half is Tomescu's: on printed
page~29 of the thesis \cite{tomescu}, which the intake read, the count of
unlabeled EADs with $n$ vertices and $s$ sources is said to vanish unless
$2^{n-s}\ge n$, equivalently $s\le n-\lceil\log_2n\rceil$, ``by
extensionality and the pigeonhole principle''. What the present proof adds
is the matching construction for every admissible source count (the
sufficiency half of Theorem~\ref{sbd:thm:support}), itself elementary; no
priority is claimed for it. The word ``conjecture'' in OEIS records the
status of an entry's formula in the encyclopedia. It does not certify that
the problem was open in the literature, nor that the proof below is the
first.

\emph{OEIS check.} On 3 October 2026 the intake read A182220 on
\url{https://oeis.org/A182220}. Its formula section lists four conjectural
expressions: Karttunen's \eqref{sbd:eq:oeis-conjecture} (2013), and others
by W.~I.~Hurt (2014), P.~Barry (2017) and A.~Krivilev (August 2026, the
``formula from August 2026'' of the package's \nolinkurl{SOURCES.md}). Each is
elementarily equivalent to \eqref{sbd:eq:max} (Barry's counts the powers of
two in $[1,n-1]$, Krivilev's the non-powers of two there), and the intake
confirmed all four against \eqref{sbd:eq:max} for $n<600$, as well as the
30 displayed terms. Theorem~\ref{sbd:thm:support} therefore settles all
four.

\emph{Repository.} The claims of this subsection still hold at the time
of writing: no other file of the repository mentions A182220, A182162,
A001192, A182161, extensional acyclic digraphs, Tomescu, Policriti or
Peddicord. No Lean or Rocq declaration formalizes any statement of this
report, and its place in the research-report collection confers no formal
status; Section~\ref{sbd:sec:formal} is a plan. Related material, none of
it a shared theorem:
(i)~Lemma~\ref{sbd:lem:collapse} is the classical Mostowski collapse, which
the surreal report
\nolinkurl{Algebra/SurrealNumbers/docs/foundations-and-computation/birthday-cutoffs-and-hereditary-sets}
proves for rooted well-founded codes as its \nolinkurl{hset:lem:collapse};
(ii)~the collection report
\nolinkurl{generating-functions-and-asymptotics/oeis-sequence-asymptotics/a116379-bounded-identity-trees}
counts rooted identity trees, the other classical coding of hereditarily
finite sets (it cites Meir, Moon and Mycielski); and
(iii)~Lemma~\ref{sbd:lem:jump} is a counterpart of a criterion of
\nolinkurl{a003407-dyadic-scaling-rigidity}, discussed in a note after it.

\emph{Notation.} Several letters carry more than one local meaning:
\begin{itemize}
\item $d$ is the defect, never the differential in $d/dp$ or $d/dz$;
 $B_d(z)$ in \eqref{sbd:eq:Bgf} is a generating function, while
 $B_{2j}$ in Theorem~\ref{sbd:thm:expansion} are Bernoulli numbers.
\item $T=2^m$ in Sections~\ref{sbd:sec:entropy} and~\ref{sbd:sec:expansion},
 but $T(n,k)$ in Appendix~\ref{sbd:app:oeis} is the labeled count
 $n!\,u(n,k)$ of A182162, as is the symbol $E(n,a(n))$ in
 \eqref{sbd:eq:extremal-count}, which the text does not define otherwise.
\item $C$ is a transitive core (and $C$, $A$ constants in the proof of
 Lemma~\ref{sbd:lem:jump}; $C_1(m,p)$ a correction coefficient);
 $L$ is a source family, $L(z)$ the lacunary series of
 Proposition~\ref{sbd:prop:maxgf}, and $L$ also a truncation order in
 Theorem~\ref{sbd:thm:expansion}.
\item $N=2^m$ in the fixed-deletion and arc sections, $N_r$ the jump
 indices of Lemma~\ref{sbd:lem:jump}, $N(0,1)$ the normal law;
 $p=n/T$ is a phase, $p_j$ the recurrence polynomials of
 \eqref{sbd:eq:prec}; $R$ a removed family, $R_d$ a radius; $M=2^m-m$
 and, in Theorem~\ref{sbd:thm:cover}, $A=u_m\binom Mk$.
\item $\rho_n\in(1/2,1]$ is the dyadic phase $n/2^{q}$, whereas $p=n/T$
 with $T=2^{q+d}$ lies in $(2^{-d-1},2^{-d}]$: the tempting reading
 ``$p=\rho_n$'' is right only for $d=0$.
\end{itemize}
No symbol was renamed.
\end{writenote}

\section{Definitions and finite-set foundations}

\subsection{Orientation conventions}
A digraph in this article is finite, loopless, and has no parallel arcs.
An arc $v\to w$ means that $w$ is an out-neighbor of $v$.
A source has indegree zero; a sink has outdegree zero.
A digraph is \emph{extensional} when the map
\[
 v\longmapsto N^+(v)=\{w:v\to w\}
\]
is injective. It is acyclic when it contains no directed cycle.
We abbreviate ``extensional acyclic digraph'' to EAD.
Every nonempty EAD has exactly one sink: a DAG has a sink, and two sinks
would have the same empty out-neighborhood. No transitive closure or
transitive reduction of the arc relation is implicit in this definition.

Let $u_m$ be the number of EAD isomorphism classes with $m$ vertices,
including $u_0=1$. We shall identify it with A001192. We use
$\binom{M}{r}=0$ if $r<0$ or $r>M$ for nonnegative integer $M$.
All logarithms without a subscript are natural.

\subsection{Collapse, transitivity, and rigidity}
\begin{lemma}[Finite collapse]\label{sbd:lem:collapse}
For every EAD $D$, the recursion
\[
 \phi(v)=\{\phi(w):w\in N^+(v)\}
\]
defines an injective map into the hereditarily finite sets. Its image
$F=\{\phi(v):v\in V(D)\}$ is a finite transitive set, and $D$ is
isomorphic to the membership digraph on $F$, with arcs $x\to y$ iff
$y\in x$. Conversely, every finite transitive set yields an EAD this way.
This gives a bijection between isomorphism classes of EADs and finite
transitive sets.
\end{lemma}
\begin{proof}
Order the vertices so that every out-neighbor precedes its parent.
The recursion is then finite and well defined. To prove injectivity,
induct on the maximum height of two vertices, where height is the
maximum length of a directed path from the vertex. Equal collapsed sets
at height at most $h$ have the same collapsed members. By the induction
hypothesis the corresponding vertices of smaller height are identical.
Thus their out-neighborhoods coincide, and extensionality identifies
the two vertices. The base case is the unique sink.

Every member of $\phi(v)$ is another element of the image, so $F$ is
transitive: $\bigcup F\subseteq F$. Injectivity shows that membership
reproduces exactly the arcs. Conversely, on a finite transitive set,
Foundation precludes membership cycles, and Extensionality says that
distinct elements have distinct member sets. These two constructions
are inverse, including at the empty set.
\end{proof}

\begin{corollary}[Rigidity]\label{sbd:cor:rigid}
Every EAD has trivial automorphism group. Consequently the number of
labeled EADs with $n$ vertices and $k$ sources is $n!u(n,k)$, and the total
number is $n!u_n$.
\end{corollary}
\begin{proof}
An automorphism preserves the recursive collapsed value of every vertex.
The collapse is injective, so the automorphism fixes every vertex.
There are therefore exactly $n!$ different labelings of each isomorphism
class, and source count is preserved by relabeling.
\end{proof}
This explains the normalization relating A182162 to $u(n,k)$ and
A182161 to A001192 \cite{oeis162,oeis161,oeis1192}.

\begin{lemma}[A sufficient bound on the small-core count]\label{sbd:lem:ubound}
For every $m\ge0$,
\begin{equation}\label{sbd:eq:smallcore-bound}
 1\le u_m\le 2^{\binom m2}.
\end{equation}
\end{lemma}
\begin{proof}
The finite ordinal $m=\{0,1,\ldots,m-1\}$ is a transitive set of size $m$,
so at least one class exists. Every DAG has a vertex ordering in which
all arcs point to earlier vertices. For a fixed such ordering there are
$2^{\binom m2}$ possible arc sets. Their isomorphism classes cover all
DAG isomorphism classes, and hence all EAD classes. Injectivity of this
covering is not needed for the upper bound.
\end{proof}

\section{The exact source maximum and all feasible source counts}

\begin{theorem}[Complete source support]\label{sbd:thm:support}
For $n\ge1$ and $k\ge1$, an EAD on $n$ vertices with exactly $k$ sources
exists if and only if
\begin{equation}\label{sbd:eq:feasible}
 k\le n,\qquad n\le 2^{n-k}.
\end{equation}
In particular,
\begin{equation}\label{sbd:eq:max}
 \boxed{\ a(n)=n-\lceil\log_2n\rceil\ }.
\end{equation}
Every integer from $1$ through $a(n)$ occurs, so A182162 has no holes
in the positive support of any row.
\end{theorem}
\begin{proof}
Suppose $D$ has $k$ sources and let $C$ be its set of nonsources,
$m=|C|=n-k$. No arc enters a source. Thus every out-neighborhood is a
subset of $C$. The injective map $v\mapsto N^+(v)$ takes $n$ vertices
into a set of size $2^m$, proving $n\le2^m$.

For sufficiency, first take $m=0$. The conditions force $n=k=1$,
and the one-vertex edgeless digraph works. Suppose $m\ge1$.
Take a directed path on $m$ core vertices, with arcs from each vertex
to the preceding one. Its $m$ out-neighborhoods are distinct. The full
set $C$ is not among them, since it would contain the vertex whose
neighborhood it was, creating a loop. There are $2^m-m$ unused subsets
of $C$, at least $k$ by the assumed inequality. Choose $k$ distinct
unused subsets, one of them equal to $C$, and add a new vertex for each
chosen subset, with exactly that subset as its out-neighborhood.

There are no arcs into the new vertices, so they are sources. All core
vertices receive an arc from the new vertex whose neighborhood is $C$,
so none is a source. Acyclicity and extensionality are immediate from
the construction. Hence the graph has exactly $k$ sources.

Finally, $n\le2^{n-k}$ is equivalent to $k\le n-q(n)$.
For $n\ge1$, this upper endpoint is at least one. This proves both the
maximum and the entire support assertion.
\end{proof}

\subsection{Elementary consequences for A182220}
For $n\ge2$, $\lceil\log_2n\rceil=1+\lfloor\log_2(n-1)\rfloor$,
which proves \eqref{sbd:eq:oeis-conjecture}. Also,
\begin{equation}\label{sbd:eq:step}
 a(n+1)-a(n)=1-\boldsymbol 1_{\{n\text{ is a power of }2\}},
 \qquad a(1)=1.
\end{equation}
This accounts for the repeated values. In particular the initial value
$1$ occurs three times, because both $1$ and $2$ are powers of two.
For an interpretation by counting nonpowers of two, zero must be included:
\[
 a(n)=\#\{0\le j\le n-1:j\text{ is not a positive power of }2\},
\]
where ``a power of $2$'' includes $2^0=1$. Omitting zero would create an
off-by-one error. We do not assert every ancillary OEIS cross-reference
without checking its conventions; the explicit formulas above remove
that ambiguity.

\subsection{A lacunary generating function for the maximum}
The next consequence is included to separate the simple maximum sequence
from the more intricate sequence counting its extremizers.

\begin{proposition}\label{sbd:prop:maxgf}
For $|z|<1$,
\begin{equation}\label{sbd:eq:maxgf}
 \sum_{n\ge1}a(n)z^n
 =\frac{z}{(1-z)^2}-\frac{z}{1-z}\sum_{j\ge0}z^{2^j}.
\end{equation}
The unit circle is a natural boundary for this generating function.
\end{proposition}
\begin{proof}
Multiply the left side by $1-z$ and apply \eqref{sbd:eq:step} to obtain the
identity. Let $L(z)=\sum_{j\ge0}z^{2^j}$ and let $\zeta$ be a root of unity
of order $2^s$. For all $j\ge s$, $\zeta^{2^j}=1$; hence
\[
 L(r\zeta)=O(1)+\sum_{j\ge s}r^{2^j},\qquad r\uparrow1,
\]
and its real part tends to infinity. At every such $\zeta\ne1$, the
rational factor $z/(1-z)$ tends to a nonzero finite value, so the right
side of \eqref{sbd:eq:maxgf} cannot extend holomorphically through $\zeta$.
These roots are dense on the circle. Analytic continuation across any
point would give a neighborhood containing one of them, a contradiction.
\end{proof}
The proof uses only a dense set of explicit singularities; no general
lacunary-series theorem is needed.

\section{Boolean-core classification and exact extremizer counts}

\subsection{The core is intrinsic}
Let $F$ be the collapsed transitive set of a graph. Its nonsources are
exactly the elements which occur as members of other elements of $F$.
Thus the nonsource core is
\begin{equation}\label{sbd:eq:union-core}
 C=\bigcup F.
\end{equation}
It is itself transitive. Indeed, $x\in y\in C$ means that $y\in z$ for
some $z\in F$; then $y\in F$, and therefore $x\in\bigcup F=C$.
If $L=F\setminus C$ is the set of sources, then
\begin{equation}\label{sbd:eq:coreform}
 F=C\cup L,\qquad L\subseteq\PP(C)\setminus C.
\end{equation}
The last exclusion uses transitivity of $C$, so that $C\subseteq\PP(C)$.
Conversely, \eqref{sbd:eq:coreform} always defines a transitive set, but its
nonsource core equals $C$ only if $\bigcup(C\cup L)=C$.
That covering condition is the sole extra constraint.

\begin{lemma}[Automatic covering above half density]\label{sbd:lem:auto}
Let $C$ be any transitive set of size $m\ge1$. If
\[
 L\subseteq\PP(C)\setminus C,\quad |L|=k,\quad m+k>2^{m-1},
\]
then $\bigcup(C\cup L)=C$.
\end{lemma}
\begin{proof}
If some $c\in C$ were outside the union, every member of the family
$C\cup L$ would avoid $c$. All $m+k$ distinct subsets would then lie in
$\PP(C\setminus\{c\})$, which contains only $2^{m-1}$ subsets.
\end{proof}

\begin{theorem}[Complete classification of extremizers]\label{sbd:thm:classify}
Let $n\ge1$, $q=\lceil\log_2n\rceil$, and $k=n-q$.
The EAD isomorphism classes attaining $a(n)$ sources are in bijection
with pairs
\begin{equation}\label{sbd:eq:extremal-pairs}
 C\in\U_q,\qquad L\in\binom{\PP(C)\setminus C}{k},
\end{equation}
where $\U_q$ is the set of transitive sets of size $q$.
Consequently
\begin{equation}\label{sbd:eq:extremal-count}
 \boxed{\ b_n^{(0)}=u_q\binom{2^q-q}{n-q},\qquad
 E(n,a(n))=n!u_q\binom{2^q-q}{n-q}\ }.
\end{equation}
\end{theorem}
\begin{proof}
Every extremizer has $q$ nonsources by Theorem~\ref{sbd:thm:support}, and
\eqref{sbd:eq:union-core}--\eqref{sbd:eq:coreform} recover its unique pair $(C,L)$.
Conversely, if $q\ge1$, then $n>2^{q-1}$, so Lemma~\ref{sbd:lem:auto}
shows that every pair in \eqref{sbd:eq:extremal-pairs} has precisely $C$ as
its core and precisely $L$ as its source set. No two different pairs give
the same transitive set, since the core is its union. Each core offers
exactly $2^q-q$ available subsets. Rigidity supplies the labeled factor.
For $n=1$, $q=0$, and the single pair is
$C=\varnothing$, $L=\{\varnothing\}$, giving the same formula.
\end{proof}

\begin{corollary}[Complete layers and sparse deletions]\label{sbd:cor:power}
For $m\ge1$ and $0\le r<2^{m-1}$,
\begin{equation}\label{sbd:eq:deletions}
 b_{2^m-r}^{(0)}=u_m\binom{2^m-m}{r}.
\end{equation}
In particular $b_{2^m}^{(0)}=u_m$, and every extremizer on $2^m$ vertices
is the membership graph of $\PP(C)$ for a unique transitive $m$-element
set $C$.
\end{corollary}
\begin{proof}
Complement the chosen source family in $\PP(C)\setminus C$.
The omitted family has size $r$, and at $r=0$ no subsets are omitted.
\end{proof}

\begin{example}
For $n=8$ we have $q=3$, $u_3=2$, and $b_8^{(0)}=2$.
For $n=9$ the core size becomes four, and
\[
 b_9^{(0)}=u_4\binom{12}{5}=9\cdot792=7128.
\]
The jump is not a numerical accident: the minimum allowed core size has
increased, making a much larger Boolean layer available. Later sections
quantify this phenomenon at every power of two.
\end{example}

\section{The marked-source sieve and finite-defect diagonals}

\subsection{A classical identity, proved in the present normalization}
\begin{proposition}[Marked sources]\label{sbd:prop:marked}
For $0\le t\le n$,
\begin{equation}\label{sbd:eq:marked}
 \sum_{s=t}^n\binom st u(n,s)
 =u_{n-t}\binom{2^{n-t}-(n-t)}{t}.
\end{equation}
For $n>0$ we set $u(n,0)=0$; for $n=0$ we set $u(0,0)=1$.
\end{proposition}
\begin{proof}
The left side counts an EAD isomorphism class together with a chosen
$t$-element subset of its sources. Rigidity ensures that different
chosen subsets cannot be identified by a nontrivial automorphism.
Delete the marked sources. Out-neighborhoods of the remaining vertices
are unchanged, since no arc entered a source, so the remaining graph is
an EAD on $n-t$ vertices.

Conversely, choose any such graph and choose $t$ distinct subsets of
its vertex set that are not already its out-neighborhoods. There are
$2^{n-t}-(n-t)$ available subsets. Adjoin one marked source for each
chosen subset. Existing sources are allowed to remain sources, which is
why no covering constraint occurs in this construction. Deleting the
marked sources reverses it. This proves the identity.
\end{proof}

The identity is the marked-source version of the classical enumeration
underlying the OEIS source sieve \cite{oeis162}. We use it as a proved
starting point rather than a new counting principle.

\begin{corollary}[Binomial inversion]\label{sbd:cor:sieve}
For $1\le k\le n$,
\begin{equation}\label{sbd:eq:sieve}
 u(n,k)=\sum_{t=k}^n(-1)^{t-k}\binom tk
 u_{n-t}\binom{2^{n-t}-(n-t)}{t}.
\end{equation}
Also
\begin{equation}\label{sbd:eq:total-rec}
 u_0=1,\qquad
 u_n=\sum_{m=0}^{n-1}(-1)^{n-m-1}
 \binom{2^m-m}{n-m}u_m\quad(n\ge1).
\end{equation}
\end{corollary}
\begin{proof}
If $P_n(z)=\sum_su(n,s)z^s$, then \eqref{sbd:eq:marked} is equivalent to
\[
 P_n(z)=\sum_{t=0}^n u_{n-t}
 \binom{2^{n-t}-(n-t)}{t}(z-1)^t.
\]
Extract $z^k$ to get \eqref{sbd:eq:sieve}. For $n>0$, evaluate at $z=0$,
use $P_n(0)=0$, and isolate the term $t=0$ to get \eqref{sbd:eq:total-rec}.
\end{proof}

\subsection{Only finitely many correction terms survive}
\begin{theorem}[Exact finite-defect formula]\label{sbd:thm:defect}
Fix $n\ge1$ and $d\ge0$ with $k=n-q-d\ge1$, and put $m=q+d$.
Then
\begin{equation}\label{sbd:eq:defect}
 \boxed{\ b_n^{(d)}=
 \sum_{j=0}^{d}(-1)^j\binom{k+j}{j}
 u_{m-j}\binom{2^{m-j}-(m-j)}{k+j}\ }.
\end{equation}
Thus an entry at distance $d$ below the source boundary requires exactly
$d+1$ displayed terms, using total counts only at the logarithmic indices
$q,q+1,\ldots,q+d$.
\end{theorem}
\begin{proof}
Substitute $t=k+j$ in \eqref{sbd:eq:sieve}. A potentially nonzero term must
satisfy
\[
 k+j\le2^{m-j}-(m-j),\quad\text{equivalently}\quad n\le2^{m-j}.
\]
Since $q=\lceil\log_2n\rceil$, this requires $m-j\ge q$, or $j\le d$.
The surviving terms are exactly those displayed. All have a positive
binomial coefficient in the stated range, though signs alternate.
\end{proof}

For $d=0$, \eqref{sbd:eq:defect} recovers the exact extremizer count, so
that count is also an immediate boundary specialization of the known
sieve. At distance one it becomes
\begin{equation}\label{sbd:eq:d1}
 b_n^{(1)}=u_{q+1}\binom{2^{q+1}-q-1}{n-q-1}
 -(n-q)b_n^{(0)}.
\end{equation}
At distance two, writing $k=n-q-2$,
\begin{align}\label{sbd:eq:d2}
 b_n^{(2)}={}&u_{q+2}\binom{2^{q+2}-q-2}{k}
 -(k+1)u_{q+1}\binom{2^{q+1}-q-1}{k+1}\notag\\
 &+\binom{k+2}{2}u_q\binom{2^q-q}{k+2}.
\end{align}
These identities should not be mistaken for asymptotic expansions:
they are exact integer formulas, even when early terms nearly cancel.

\section{Uniform covering estimates and random cores}

\subsection{An estimate valid away from the extremal boundary as well}
\begin{theorem}[Uniform core-covering bound]\label{sbd:thm:cover}
Let $m\ge0$, $k\ge1$, $n=m+k$, and $k\le M=2^m-m$.
Set $A=u_m\binom Mk$. Then
\begin{equation}\label{sbd:eq:cover}
 \boxed{\quad 0\le A-u(n,k)\le m\,2^{-k}A.\quad}
\end{equation}
For $m\ge1$, a sharper upper bound on the relative deficit is
\begin{equation}\label{sbd:eq:cover-sharp}
 \frac{A-u(n,k)}{A}
 \le\frac{\binom{2^{m-1}-m}{k}}{\binom{2^m-m}{k}}
 \frac{\sum_s s\,u(m,s)}{u_m}.
\end{equation}
\end{theorem}
\begin{proof}
For a fixed transitive core $C$ of size $m$, choose a uniform $k$-subset
$L$ of $\PP(C)\setminus C$. The resulting graph is counted by $u(n,k)$
exactly when $\bigcup(C\cup L)=C$.
Every element already in $\bigcup C$ is covered. If
$c\in C\setminus\bigcup C$, then none of the $m$ members of $C$ contains
$c$. Of the $2^{m-1}$ subsets avoiding $c$, precisely $m$ are therefore
forbidden core subsets. The probability that no chosen subset contains
$c$ is
\[
 \frac{\binom{2^{m-1}-m}{k}}{\binom{2^m-m}{k}}\le2^{-k}.
\]
If the numerator vanishes the assertion is immediate. Otherwise compare
successive factors of the two falling factorials; each ratio is at most
$1/2$. Apply the union bound over the sources of $C$, whose number is
at most $m$, and sum over all $u_m$ cores. The sharper average-source
version follows before replacing the number of sources by $m$.
For $m=0$, the only admissible choice is $k=1$, and equality holds.
\end{proof}

The average core-source count in \eqref{sbd:eq:cover-sharp} is itself exact:
Proposition~\ref{sbd:prop:marked}, with $t=1$, gives
\begin{equation}\label{sbd:eq:mean-sources}
 \frac{\sum_s s\,u(m,s)}{u_m}
 =\frac{(2^{m-1}-m+1)u_{m-1}}{u_m}\qquad(m\ge1).
\end{equation}
In applications where $k-\log_2m\to+\infty$, the simpler bound already
proves
\begin{equation}\label{sbd:eq:general-equivalent}
 u(m+k,k)=u_m\binom{2^m-m}{k}
 \left(1+O(m2^{-k})\right).
\end{equation}
The implicit error here is one-sided: the true value never exceeds the
main expression.

\subsection{Uniform sampling and the core law}
There is an exact rejection sampler for a uniformly random isomorphism
class with $n=m+k$ vertices and $k$ sources. First draw a uniform core
from $\U_m$, then a uniform $k$-subset of its $M$ available neighborhoods.
Keep the pair if it covers the core; otherwise repeat. Before rejection,
every pair has exactly the same probability, and every accepted graph
corresponds to one pair. Thus the accepted graph is exactly uniform.
If $m2^{-k}<1$, the expected number of trials is at most
$(1-m2^{-k})^{-1}$. This is conditional on having a procedure for uniformly
sampling the much smaller core; it is not a claim of a complete
polynomial-time sampler for arbitrary $m$.

At the extremal boundary, rejection is never necessary, and the core
of a uniform extremizer is \emph{exactly} uniform on $\U_q$.
More generally, write $\delta$ for the actual rejection probability.
Conditioning a finite probability distribution on an event of probability
$1-\delta$ changes it in total variation by exactly $\delta$.
Taking the core marginal cannot increase total variation. Therefore
\begin{equation}\label{sbd:eq:core-tv}
 \TV\bigl(\text{core law of a uniform }(n,k)\text{ EAD},
               \operatorname{Unif}(\U_m)\bigr)
 \le m2^{-k}.
\end{equation}
This bound is uniform over all core structures. It explains why small-core
statistics can be transferred to the high-source part of the triangle.

\section{The dyadic entropy law}\label{sbd:sec:entropy}

\subsection{Uniform logarithmic asymptotics}
Define the binary entropy, in natural units, by
\[
 H(p)=-p\log p-(1-p)\log(1-p),\qquad H(0)=H(1)=0.
\]
For a fixed defect $d$, set
\begin{equation}\label{sbd:eq:phidef}
 \rho_n=\frac{n}{2^{q(n)}}\in(1/2,1],\qquad
 \Phi_d(\rho)=\frac{2^d}{\rho}H\!\left(\frac{\rho}{2^d}\right).
\end{equation}
The strict lower endpoint for $\rho_n$ does not preclude subsequences
converging to $1/2$.

\begin{theorem}[Entropy profile]\label{sbd:thm:entropy}
For every fixed $d\ge0$, as $n\to\infty$,
\begin{equation}\label{sbd:eq:entropy}
 \boxed{\quad \log b_n^{(d)}
       =n\Phi_d(\rho_n)+O_d((\log n)^2).\quad}
\end{equation}
The bound is uniform over the entire dyadic block, including its
endpoints. No asymptotic formula for $u_m$ is needed.
\end{theorem}
\begin{proof}
Let $m=q+d$, $T=2^m$, and $k=n-m$. These parameters satisfy
$m=O_d(\log n)$, $n\le T<2^{d+1}n$, and $k=n-O_d(\log n)$.
Theorem~\ref{sbd:thm:cover} gives
\begin{equation}\label{sbd:eq:loglead}
 \log b_n^{(d)}=\log u_m+\log\binom{T-m}{n-m}
                  +O_d(m2^{-k}).
\end{equation}
The small-core logarithm is $O_d((\log n)^2)$ by
Lemma~\ref{sbd:lem:ubound}. The exact identity
\begin{equation}\label{sbd:eq:binratio}
 \binom{T-m}{n-m}=\binom Tn\frac{\fall nm}{\fall Tm}
\end{equation}
shows that the two binomial logarithms differ by $O_d(m)$:
every factor $(n-i)/(T-i)$ is bounded below by a positive constant
depending only on $d$, for all sufficiently large $n$.

For completeness, the elementary entropy bounds are
\begin{equation}\label{sbd:eq:entropy-binom}
 \frac{e^{T H(n/T)}}{T+1}\le\binom Tn\le e^{T H(n/T)}.
\end{equation}
For $0<n<T$, consider a $\Bin(T,n/T)$ random variable. Its probability
at $n$ is a maximum probability, hence is at least $1/(T+1)$ and at most
one. Rearranging gives \eqref{sbd:eq:entropy-binom}; the endpoint cases are
immediate. It follows that $\log\binom Tn=T H(n/T)+O(\log T)$.
Finally $T H(n/T)=n\Phi_d(\rho_n)$. Combining these estimates proves
\eqref{sbd:eq:entropy}, with no loss of uniformity as $n/T\to1$.
\end{proof}

\subsection{Every subsequential exponential rate}
Set
\begin{equation}\label{sbd:eq:lambdadef}
 \lambda_d=2^d H(2^{-d}),\qquad d\ge0.
\end{equation}
In particular $\lambda_0=0$, $\lambda_1=\log4$, and
$e^{\lambda_2}=256/27$.

\begin{theorem}[Complete growth-rate interval]\label{sbd:thm:rates}
For each fixed $d\ge0$, the set of subsequential limits of
$(b_n^{(d)})^{1/n}$ is exactly
\begin{equation}\label{sbd:eq:rates}
 [e^{\lambda_d},e^{\lambda_{d+1}}].
\end{equation}
Thus
\begin{equation}\label{sbd:eq:extreme-rates}
 \liminf_{n\to\infty}(b_n^{(0)})^{1/n}=1,
 \qquad
 \limsup_{n\to\infty}(b_n^{(0)})^{1/n}=4.
\end{equation}
No fixed-defect boundary sequence has a single exponential growth
constant.
\end{theorem}
\begin{proof}
Write $p=\rho/2^d$. Then $\Phi_d(\rho)=H(p)/p$.
For $0<p<1$,
\[
 \frac{d}{dp}\frac{H(p)}p
 =\frac{pH'(p)-H(p)}{p^2}
 =\frac{\log(1-p)}{p^2}<0.
\]
Hence $\Phi_d$ is strictly decreasing, with values $\lambda_{d+1}$ at
$\rho=1/2$ and $\lambda_d$ at $\rho=1$.
Every $\rho\in(1/2,1)$ is the limit of phases obtained by rounding
$\rho2^q$ to integers. Powers $n=2^q$ realize phase one, while
$n=2^{q-1}+1$ tends to phase $1/2$. The entropy theorem gives the entire
claimed interval. Conversely, compactness of $[1/2,1]$ and continuity of
$\Phi_d$ exclude any other subsequential limits.
\end{proof}

In particular the first three intervals, for $d=0,1,2$, are
\[
 [1,4],\qquad[4,256/27],\qquad
 [256/27,\,2^{24}/7^7].
\]
The upper endpoint for one defect equals the lower endpoint for the next.
This expresses the fact that crossing a dyadic boundary adds one core
vertex, just as increasing the allowed source deficit does.

\section{Non-P-recursiveness of every fixed-defect boundary}

A sequence $(c_n)$ is P-recursive if, for all sufficiently large $n$,
\begin{equation}\label{sbd:eq:prec}
 \sum_{j=0}^{h}p_j(n)c_{n+j}=0
\end{equation}
for polynomials $p_j$, not all zero. Allowing a finite exceptional initial
segment makes the conclusion below stronger, not weaker.
A formal power series is D-finite if it satisfies a nonzero homogeneous
linear differential equation with polynomial coefficients.

\begin{lemma}[An exponential-jump obstruction]\label{sbd:lem:jump}
Suppose a nonnegative sequence $(c_n)$ has integers $N_r\to\infty$ and
constants $\alpha<\beta$ such that, for every fixed $j\ge0$,
\[
 c_{N_r-j}=\exp(\alpha N_r+o(N_r)),\qquad
 c_{N_r+1}=\exp(\beta N_r+o(N_r)).
\]
Then $(c_n)$ is not P-recursive.
\end{lemma}
\begin{proof}
Suppose \eqref{sbd:eq:prec} holds, and choose $h$ as the largest index with
$p_h\ne0$. If $h=0$, positivity along $(N_r)$ contradicts the recurrence.
Otherwise apply it at $n=N_r+1-h$. For sufficiently large $r$, $p_h(n)$
is nonzero, and ratios of fixed polynomials give a bound
\[
 c_{N_r+1}\le C N_r^A\sum_{j=0}^{h-1}c_{N_r+1-h+j}
             =\exp(\alpha N_r+o(N_r))
\]
for constants $C,A$. This contradicts $\beta>\alpha$.
Absolute values of the coefficients were used, so cancellation cannot
invalidate the estimate. The argument also applies to complex
polynomial coefficients.
\end{proof}

\begin{writenote}[Added 3 October 2026, batch~85]
The collection report
\nolinkurl{generating-functions-and-asymptotics/oeis-sequence-asymptotics/a003407-dyadic-scaling-rigidity}
proves, in its Part~II, an obstruction that runs the other way. Its
Corollary \nolinkurl{dsr:nh:cor:jumps} (a consequence of its local-growth
rigidity theorem \nolinkurl{dsr:nh:thm:rigidity}, which rests on a G-function
input) says that an eventually positive, exponentially bounded, integer
P-recursive sequence without an $n$th-root limit must have
$u_{n+1}/u_n>e^{\varepsilon n}$ for infinitely many $n$: upward
exponential jumps are \emph{necessary} for such oscillation, and their
absence rules P-recursiveness out. Lemma~\ref{sbd:lem:jump} is a
\emph{sufficient} condition in the opposite situation: a jump preceded by a
plateau of every fixed length at the lower rate rules P-recursiveness out.
It needs neither integrality nor an exponential bound, and its proof is
elementary. The two criteria are complementary, not competing. The
boundary sequences $b^{(d)}$ do have the upward jumps (from
$e^{\lambda_dN_r}$ to $e^{\lambda_{d+1}N_r}$ at $N_r=2^r$), so the a003407
criterion cannot exclude their P-recursiveness, and
Lemma~\ref{sbd:lem:jump} is the one that applies. Conversely, the
3AP-free permutation counts of a003407 have only subexponential forward
increases, so Lemma~\ref{sbd:lem:jump} does not apply to them. The
plateau hypothesis ``for every fixed $j\ge0$'' is essential: a003407's
example $3^n+(-3)^n+2^n$, which satisfies a constant-coefficient
recurrence, jumps from rate $\log2$ at odd indices to $\log3$ at even
ones, but its two-term window never stays at the lower rate.
\end{writenote}

\begin{theorem}[Non-holonomicity]\label{sbd:thm:nonholonomic}
For every fixed $d\ge0$, neither $(b_n^{(d)})$ nor $(n!b_n^{(d)})$ is
P-recursive. The ordinary generating function
\begin{equation}\label{sbd:eq:Bgf}
 B_d(z)=\sum_{n\ge1}b_n^{(d)}z^n
\end{equation}
is not D-finite. Its radius of convergence is
\begin{equation}\label{sbd:eq:radius}
 R_d=e^{-\lambda_{d+1}};
 \qquad R_0=\frac14.
\end{equation}
\end{theorem}
\begin{proof}
Take $N_r=2^r$. For every fixed $j\ge0$, the index $N_r-j$ remains in
the dyadic block with $q=r$ when $r$ is large, and its phase tends to
one. By Theorem~\ref{sbd:thm:entropy},
\[
 \log b_{N_r-j}^{(d)}=\lambda_d N_r+o(N_r).
\]
At $N_r+1$, $q=r+1$ and the phase tends to $1/2$, giving
\[
 \log b_{N_r+1}^{(d)}=\lambda_{d+1}N_r+o(N_r).
\]
The strict decrease proved in Theorem~\ref{sbd:thm:rates} yields
$\lambda_{d+1}>\lambda_d$. Lemma~\ref{sbd:lem:jump} applies.

If $n!b_n^{(d)}$ satisfied a polynomial recurrence, division by $n!$
would give one for $b_n^{(d)}$: each factor $(n+j)!/n!$ is a polynomial
in $n$. Thus the labeled sequence is also non-P-recursive.

For completeness, extracting coefficients from a nonzero polynomial
linear differential equation for $B_d$ yields a nonzero finite-shift
polynomial recurrence for its coefficients. Indeed, coefficients of
$z^i(d/dz)^jB_d$ are shifted coefficients multiplied by falling
factorials in the index; a common shift removes any negative indices.
This contradicts what was just proved. Finally, the Cauchy--Hadamard
formula and Theorem~\ref{sbd:thm:rates} give \eqref{sbd:eq:radius}.
\end{proof}

\begin{remark}[What the theorem does not say]
Non-D-finiteness does not by itself imply that the circle $|z|=R_d$ is a
natural boundary, nor does it rule out useful nonlinear functional
equations. We prove a natural boundary for the maximum sequence in
Proposition~\ref{sbd:prop:maxgf}, but do not make that stronger claim for
$B_d$. The labeled ordinary generating series has radius zero; the
series $B_d$ is also the exponential generating series of the labeled
boundary counts.
\end{remark}

\section{An all-orders expansion inside a dyadic block}\label{sbd:sec:expansion}

The entropy law is uniform even at complete layers. More detailed
prefactors require staying away from the degeneracy $p=1$.
This section retains $u_m$ exactly; it does not purport to replace the
classical, separate asymptotic analysis of the total EAD count.

\begin{theorem}[Explicit expansion with exact core amplitude]\label{sbd:thm:expansion}
Fix $d\ge0$ and $\varepsilon>0$. Let
\[
 m=q+d,\quad T=2^m,\quad p=n/T,\quad k=n-m,
 \qquad \varepsilon\le p\le1-\varepsilon.
\]
Define $S_\ell(m)=\sum_{i=0}^{m-1}i^\ell$ and let $B_{2j}$ denote the
Bernoulli numbers with $B_2=1/6$. For fixed integers $L\ge0$, $J\ge1$,
\begin{align}\label{sbd:eq:allorders}
 \log b_n^{(d)}={}&\log u_m+T H(p)
 -\tfrac12\log\bigl(2\pi T p(1-p)\bigr)+m\log p\notag\\
 &-\sum_{\ell=1}^{L}
       \frac{p^{-\ell}-1}{\ell T^\ell}S_\ell(m)\notag\\
 &+\sum_{j=1}^{J}\frac{B_{2j}}{2j(2j-1)T^{2j-1}}
       \bigl[1-p^{1-2j}-(1-p)^{1-2j}\bigr]\notag\\
 &+O_{d,\varepsilon,L,J}\left(
       \frac{m^{L+2}}{T^{L+1}}+T^{-(2J+1)}+m2^{-k}\right).
\end{align}
The first sum is empty when $L=0$.
\end{theorem}
\begin{proof}
Start with \eqref{sbd:eq:loglead} and \eqref{sbd:eq:binratio}.
The logarithm of the falling-factorial ratio is
\[
 m\log p+\sum_{i=0}^{m-1}
 \left[\log\left(1-\frac{i}{pT}\right)
             -\log\left(1-\frac{i}{T}\right)\right].
\]
Since $m=O_d(\log T)$ and $p\ge\varepsilon$, expand both logarithms
through degree $L$. The coefficient of $T^{-\ell}$ is exactly the
first sum in \eqref{sbd:eq:allorders}. Summing the elementary Taylor
remainders gives $O(m^{L+2}/T^{L+1})$.

Apply the Stirling expansion for $\log(T!)$, $\log((pT)!)$, and
$\log(((1-p)T)!)$. Subtraction produces $T H(p)$, the displayed
Gaussian prefactor, and the Bernoulli sum. Its remainder after $J$
terms is $O_{\varepsilon,J}(T^{-(2J+1)})$. This standard expansion follows,
for example, by Euler--Maclaurin summation of $\log x$; its remainder
has the asserted bound because $pT$ and $(1-p)T$ are both at least
$\varepsilon T$. Finally include the covering error from
\eqref{sbd:eq:loglead}.
\end{proof}

In particular,
\begin{align}\label{sbd:eq:first-prefactor}
 b_n^{(d)}={}&
 \frac{u_m p^m e^{T H(p)}}{\sqrt{2\pi T p(1-p)}}
 \left[1+\frac{C_1(m,p)}T
             +O_{d,\varepsilon}\left(\frac{m^4}{T^2}\right)\right],\\
 C_1(m,p)={}&-\frac{1-p}{2p}m(m-1)
       +\frac{1-p^{-1}-(1-p)^{-1}}{12}.\notag
\end{align}
Taking larger $L$ and $J$ supplies arbitrarily many logarithmic
corrections. Exponentiating gives a multiplicative expansion whose
coefficients are explicit polynomials in $m$, with rational functions
of $p$ as coefficients.

For $d\ge1$, the actual phases have
$p\in(2^{-d-1},2^{-d}]$, so a suitable fixed $\varepsilon$ covers the
entire dyadic block. For $d=0$, the complete-layer edge $p\to1$ must
instead be treated by the exact deletion formula
\eqref{sbd:eq:deletions}. This is a genuine boundary change of asymptotic
regime, not an error to be removed by formally substituting $p=1$ in
the Gaussian prefactor.

\subsection{The fixed-deletion regime}
For fixed $r\ge0$ and $N=2^m$, Corollary~\ref{sbd:cor:power} gives
\begin{equation}\label{sbd:eq:fixedr}
 b_{N-r}^{(0)}
 =u_m\frac{N^r}{r!}
 \left[1-\frac{rm+r(r-1)/2}{N}
          +O_r\left(\frac{m^2}{N^2}\right)\right].
\end{equation}
This follows directly by expanding
$\binom{N-m}{r}=N^r/r!\prod_{i=0}^{r-1}(1-(m+i)/N)$.
More generally, for $r=o(N)$ with $m+r<N$, an exact logarithmic series is
\begin{equation}\label{sbd:eq:deletion-series}
 \log\frac{r!\,b_{N-r}^{(0)}}{u_mN^r}
 =-\sum_{\ell\ge1}\frac{1}{\ell N^\ell}
       \sum_{i=0}^{r-1}(m+i)^\ell.
\end{equation}
It converges because $(m+r-1)/N<1$, and truncation is controlled by the
geometric tail of the logarithm series. This supplies a separate exact
route toward a matched sparse-to-bulk expansion.

\section{Arc counts and a binomial law near complete layers}

\subsection{A core-dependent polynomial with a simple complement form}
Let $e(F)$ be the number of arcs in the membership graph of $F$.
Since arcs are membership incidences,
\[
 e(F)=\sum_{S\in F}|S|.
\]
For $n=2^m-r$ with $0\le r<2^{m-1}$, an extremizer is
$F=\PP(C)\setminus R$, where $C\in\U_m$ and
$R\subseteq\PP(C)\setminus C$ has size $r$. Hence
\begin{equation}\label{sbd:eq:arc-deficit}
 e(F)=m2^{m-1}-Z,\qquad Z=\sum_{S\in R}|S|.
\end{equation}
In particular, \emph{all} extremizers at $n=2^m$ have exactly
$m2^{m-1}$ arcs, regardless of the core graph.

For a fixed core $C$, the full deficit-counting polynomial is
\begin{equation}\label{sbd:eq:arc-poly}
 [z^r]\prod_{S\in\PP(C)\setminus C}(1+z w^{|S|}).
\end{equation}
Summing this over $C\in\U_m$ gives the exact unlabeled joint count by
size and arc deficit. Equivalently, the product is the formal quotient
\[
 \frac{\displaystyle\prod_{j=0}^{m}(1+z w^j)^{\binom mj}}
      {\displaystyle\prod_{x\in C}(1+z w^{|x|})}.
\]
The denominator divides the numerator as a polynomial product; it is
not an analytic singularity assertion.

\begin{theorem}[Total-variation binomial approximation]\label{sbd:thm:binomial}
Let $F$ be uniform among extremal EAD isomorphism classes on $n=2^m-r$
vertices, where $m\ge1$ and $0\le r<2^{m-1}$. Define
$Z=m2^{m-1}-e(F)$. Then
\begin{equation}\label{sbd:eq:tvbin}
 \boxed{\quad
 \TV\bigl(\mathcal L(Z),\Bin(rm,1/2)\bigr)
 \le \min\left\{1,\frac{rm+\binom r2}{2^m}\right\}.
 \quad}
\end{equation}
The same statement holds for a uniformly labeled extremizer.
\end{theorem}
\begin{proof}
By Theorem~\ref{sbd:thm:classify}, the core is uniform on $\U_m$, and,
conditional on it, $R$ is a uniform $r$-subset of the available family.
Fix a core and draw $r$ independent uniform subsets
$S_1,\ldots,S_r$ of the entire Boolean cube $\PP(C)$.
Each $S_i$ has a $\Bin(m,1/2)$ size, independently; therefore
$\sum_i|S_i|$ is exactly $\Bin(rm,1/2)$.

Condition this ordered sample on avoiding the $m$ forbidden core
subsets and on having no repetition. The conditional distribution is
exactly uniform sampling without replacement from
$\PP(C)\setminus C$, with an irrelevant ordering. By the union bound,
the conditioning event fails with probability at most
\[
 r\frac{m}{2^m}+\binom r2\frac1{2^m}.
\]
Conditioning changes total variation by the failure probability, and
mapping to the sum of the sizes cannot increase total variation.
The estimate is independent of the core, so it survives mixing over
cores. Rigidity implies that uniform labeling gives the same induced
law on isomorphism classes, proving the last assertion.
\end{proof}

\begin{corollary}[Boundary central limit law]\label{sbd:cor:clt}
If $m\to\infty$, $r=r(m)\ge1$, and $r(m+r)=o(2^m)$, then
\begin{equation}\label{sbd:eq:clt}
 \frac{m2^{m-1}-e(F)-rm/2}{\sqrt{rm}/2}
 \ \Longrightarrow\ N(0,1).
\end{equation}
\end{corollary}
\begin{proof}
The total-variation error in \eqref{sbd:eq:tvbin} tends to zero.
For $B\sim\Bin(rm,1/2)$, the characteristic function of
$(B-rm/2)/(\sqrt{rm}/2)$ is
$\cos(t/\sqrt{rm})^{rm}$, which tends to $e^{-t^2/2}$.
The claimed convergence follows, since $rm\to\infty$.
\end{proof}
For fixed $r\ge1$, the error is $O_r(m2^{-m})$ and the Gaussian limit
already applies. For $r=0$, the distribution is instead deterministic.
The theorem concerns graphs conditioned to maximize their source count;
it is not Wagner's normal limit for unrestricted EADs.

\subsection{Exact conditional moments}
Let $N=2^m$, $M=N-m$, and define two core statistics
\[
 e_C=\sum_{x\in C}|x|,\qquad h_C=\sum_{x\in C}|x|^2.
\]
The sum and squared sum of subset sizes in the available population are
\[
 A_1=\frac{mN}{2}-e_C,\qquad
 A_2=\frac{(m^2+m)N}{4}-h_C.
\]
Writing $\mu_C=A_1/M$ and $v_C=A_2/M-\mu_C^2$, uniform
sampling without replacement gives
\begin{equation}\label{sbd:eq:moments}
 \E(Z\mid C)=r\mu_C,\qquad
 \Var(Z\mid C)=\frac{r(M-r)}{M-1}v_C\quad(M>1).
\end{equation}
The variance identity follows by expanding the variance of the sum of
$r$ sampled population entries: the covariance of two distinct sampled
entries is $-v_C/(M-1)$. This derivation also specifies the variance
normalization. When $M=1$, the sample is deterministic.
These exact moments remain available outside the small-error range of
Theorem~\ref{sbd:thm:binomial}.

\section{Exact verification and computational artifacts}

\subsection{Three complementary checks}
The accompanying program performs the following distinct checks.
First, it computes $u_n$ from \eqref{sbd:eq:total-rec} and then evaluates the
marked-source sieve. Second, it computes the source triangle using the
source-deletion recurrence from Tomescu's Corollary 2.1.7
\cite{tomescu}:
\begin{align}\label{sbd:eq:deletionrec}
 u(n,s)=\frac1s\bigg[&\bigl(2^{n-s}-(n-1)\bigr)u(n-1,s-1)\notag\\
 &+\sum_{j=0}^{n-s-1}\binom{s+j}{j+1}
              2^{n-1-s-j}u(n-1,s+j)\bigg].
\end{align}
Here $u(1,1)=1$ and out-of-range entries are zero. The code verifies
exact divisibility by $s$, equality of the two triangle constructions,
the complete support theorem, the finite-defect formulas, and the
covering inequalities.

Third, for every $n\le7$, the program enumerates all graphs in which
arcs point from a larger vertex index to a smaller one. It rejects
repeated out-neighborhoods, collapses surviving graphs to nested finite
sets, and deduplicates the resulting isomorphism classes. This does not
use either counting recurrence. Source counts and arc counts are
computed directly from the arc masks. Every DAG has such an ordering,
so the enumeration covers every class, even though some classes occur
under several orderings.

\begin{table}[htbp]
\centering
\caption{Exhaustive enumeration, compared with the recurrence counts.
The middle column counts surviving graphs in a fixed topological
labeling convention, before isomorphism deduplication.}
\begin{tabular}{r r r l}
\toprule
$n$ & Surviving ordered graphs & Classes $u_n$ & Positive source row $u(n,k)$\\
\midrule
1&1&1&1\\
2&1&1&1\\
3&2&2&2\\
4&10&9&8, 1\\
5&120&88&68, 20\\
6&3240&1802&1248, 534, 20\\
7&187920&75598&48640, 25008, 1940, 10\\
\bottomrule
\end{tabular}
\end{table}

\subsection{Recorded results}
The recorded run used \code{--max-n 64 --brute-n 7}. It checked
2080 triangle cells by both recurrences, 2080 covering inequalities,
64 extremizer counts, and 354 finite-defect entries with $d\le5$.
It also matched the first 17 displayed terms of A001192 and the first
25 displayed flattened labeled entries of A182162
\cite{oeis1192,oeis162}.
For each $n\le7$, the directly enumerated extremal arc polynomial
agreed with the Boolean-core construction.

\begin{table}[htbp]
\centering
\caption{Exact unlabeled extremizer counts. Labeled counts are obtained
by multiplying the final column by $n!$.}
\begin{tabular}{r r r r@{\qquad}r r r r}
\toprule
$n$&$q$&$a(n)$&$b_n^{(0)}$&$n$&$q$&$a(n)$&$b_n^{(0)}$\\
\midrule
1&0&1&1&9&4&5&7128\\
2&1&1&1&10&4&6&8316\\
3&2&1&2&11&4&7&7128\\
4&2&2&1&12&4&8&4455\\
5&3&2&20&13&4&9&1980\\
6&3&3&20&14&4&10&594\\
7&3&4&10&15&4&11&108\\
8&3&5&2&16&4&12&9\\
\bottomrule
\end{tabular}
\end{table}

The source-count CSV files contain exact integers, not rounded plots or
floating-point approximations. A separate \code{check_asymptotics.py}
uses 70-digit Decimal arithmetic to test the first correction at
$n=192,768,3072,12288$ and $d=0,1,2$; all twelve tests improve the
leading approximation. Its separate CSV records numerical diagnostics,
not interval-certified bounds. The JSON report records the checks and
their ranges. No finite experiment is presented as a proof of
non-P-recursiveness, the entropy theorem, or a limit distribution.

\subsection{Computational complexity and numerical stability}
For fixed $d$, \eqref{sbd:eq:defect} uses only $d+1$ terms after the total
counts through $q+d$ have been tabulated. This is an arithmetic-operation
statement, not a bit-complexity claim: the binomial coefficients and the
answers themselves have $\Theta_d(n)$ bits in the bulk. Near a complete
layer, the exact complement form \eqref{sbd:eq:deletions} can be much smaller.

Use integer binomial coefficients for exact evaluation. For large
floating-point evaluations, use logarithms and either the covering
approximation or the all-orders expansion; do not subtract enormous,
nearly equal floating-point source-sieve terms near small parameters.
The program caps brute-force enumeration at seven vertices to make the
exponential nature of that check explicit rather than silently attempting
an infeasible extension.

\begin{writenote}[Added 3 October 2026, batch~85]
In the repository the two programs ship as \nolinkurl{code/verify.py} and
\nolinkurl{code/check_asymptotics.py}, and their outputs under \nolinkurl{data/}.
The delivered \nolinkurl{data/verification-output.txt} is not shipped: it was
a byte copy of \nolinkurl{data/verification.json}, since the program prints the
JSON report it writes. Both programs write beside themselves (into a
\nolinkurl{data} directory next to the script, the second without an option),
so the report's README gives a rerun on a scratch copy in the delivered
layout. The intake ran them so on Python~3.14.4: \code{verify.py --max-n 64 --brute-n 7}
passed in about 7\,s, \nolinkurl{check_asymptotics.py} in
about 4\,s; the three CSV files were reproduced byte for byte, the JSON
report up to Windows line endings, and both printed transcripts matched the
recorded ones. The 17 terms of A001192 and 25 terms of A182162 that
\nolinkurl{verify.py} embeds as test fixtures are OEIS data, published under
CC~BY-SA~4.0.
\end{writenote}

\section{Formalization route and further research}

\subsection{A finite, auditable formalization route}\label{sbd:sec:formal}
A formal development need not begin with the full set-theoretic collapse.
Represent a finite digraph by a Boolean relation on a finite vertex type.
Define the finite set of out-neighbors, sourcehood, acyclicity, and
extensionality. The maximum-source proof then reduces to an injection
from vertices into the finite powerset of the nonsources, followed by
the explicit construction in Theorem~\ref{sbd:thm:support}.

A useful next layer would prove that deletion of sources preserves
out-neighborhoods and extensionality. This supports a marked-source
bijection and finite binomial inversion. To treat unlabeled counts,
one may either prove rigidity and use labeled counts throughout, or
introduce finite collapse and quotient by graph isomorphism. The former
avoids set-theoretic infrastructure for the first main counting results.

The entropy and non-holonomicity proofs can be separated from graph
semantics. A generic exponential-jump lemma is particularly suitable for
reuse: its hypotheses are scalar asymptotic statements, and its proof
only needs eventual polynomial-ratio bounds and the triangle inequality.
The deliverable contains no Lean source and makes no claim of kernel
verification. This paragraph is a development plan, not an assertion that
particular library declarations already implement the argument.

\subsection{Specific research questions}\label{sbd:sec:questions}
The following questions are proposed directions, not claims that each
has been certified open in the published literature.

\paragraph{1. Natural boundaries of the extremizer generating functions.}
Does $B_d(z)$ have a natural boundary on $|z|=e^{-\lambda_{d+1}}$?
The proof here establishes non-D-finiteness but not a dense set of
singularities. Unlike the lacunary maximum sequence, $B_d$ has positive
coefficients throughout its eventual support. A useful first target is
an analysis of contributions from short intervals immediately after
powers of two.

\paragraph{2. A sharper small-core amplitude.}
Combine an explicit, verified asymptotic expansion for $u_m$ with
\eqref{sbd:eq:allorders}. Since $m\asymp\log n$, the resulting corrections
live on scales different from the Boolean-layer saddle. Distinguish
statements uniform in the dyadic phase from fixed-phase expansions.

\paragraph{3. Defects growing with the graph size.}
How far can the entropy law and exponential-jump obstruction be extended
when $d=d(n)$ increases? The covering theorem itself remains valid, but
$\log u_{q+d}$ ceases to be a negligible amplitude if the core becomes
too large. Identify the transition at which the large-source boundary
analysis must connect with unrestricted EAD enumeration.

\paragraph{4. A matched bulk-to-complete-layer expansion.}
The bulk formula and the fixed-deletion formula overlap when the number
$r=2^m-n$ tends to infinity slowly. Derive explicit uniform remainder
bounds covering $r$ from fixed values to a positive fraction of $2^m$,
without using a Gaussian prefactor at a degenerate endpoint.

\paragraph{5. Extend the arc law beyond the collision estimate.}
The total-variation coupling loses strength when $r$ is comparable to
$2^{m/2}$. This is a limitation of comparing whole subset samples, not
necessarily of comparing their sizes. Can a finite-population local
limit theorem recover the arc distribution for all densities, with an
explicit correction for the $m$ forbidden core subsets?

\paragraph{6. Joint degree statistics.}
Study the empirical outdegree profile of source vertices and the joint
indegree distribution of core vertices. The source-neighborhood family
is a sample without replacement from a Boolean cube with a small deleted
family. This suggests multivariate hypergeometric limits, but the exact
core-induced correlations should be retained rather than discarded.

\paragraph{7. Source layers under repeated deletion.}
Successively delete all sources. Conditional on the first layer being
extremal, the remaining graph is exactly a uniform small core. Can this
be iterated to obtain precise distributions for the number of rounds,
the full layer-size vector, or the maximum rank under extremal source
conditioning?

\paragraph{8. Restricted degree and weak extensionality.}
If outdegrees are at most $\Delta$, the Boolean capacity becomes
$\sum_{j\le\Delta}\binom mj$ rather than $2^m$. Determine when its
pigeonhole bound is sharp, how covering changes, and whether analogous
boundary jumps imply non-P-recursiveness. A related problem permits
bounded multiplicity of out-neighborhoods or distinguished atoms.

\paragraph{9. Exact ranking and sampling of boundary classes.}
Implement ranking and unranking for the core-plus-subset representation.
A ranked core together with a combinatorial-number-system rank for its
chosen neighborhoods gives a natural interface. Specify bit complexity,
including the cost of representing deeply nested finite sets.

\paragraph{10. Formal proofs and OEIS editorial updates.}
Formalize the source support theorem and the boundary bijection first.
After independent mathematical review, the proof of A182220 and the
boundary formula can be proposed to OEIS with the historical upper-bound
citation. No OEIS submission or repository commit was made during this
work.

\section{Conclusion}
The source maximum has a short Boolean-capacity proof, and its known
upper bound should not be confused with a new breakthrough. The
interesting structure appears when one counts the graphs attaining,
or nearly attaining, that maximum. Their core has logarithmic size,
while their sources select a large family of subsets of that core.
At maximum density, coverage is automatic and the count is exactly
$u_q\binom{2^q-q}{n-q}$. At fixed defect, only finitely many terms of the
classical sieve survive.

This decomposition makes it possible to prove a uniform entropy law
without solving the separate asymptotic problem for the core count.
Its phase resets at powers of two. The resulting exponential jumps
rule out every polynomial-coefficient linear recurrence for each fixed
boundary diagonal. Exact arc formulas and the binomial coupling provide
further information about the extremizers themselves.

The principal deliverable is therefore a chain of proved results:
sharp source support, exact boundary enumeration, finite-defect
formulas, uniform high-source estimates, complete growth-rate intervals,
and non-holonomicity, with reproducible finite checks. Questions about
historical priority, natural boundaries of the extremizer series,
growing defects, and full formal verification remain explicitly separate.

\appendix
\section{Suggested concise OEIS-facing statements}\label{sbd:app:oeis}

\paragraph{A182220.}
For every $n\ge1$, $a(n)=n-\lceil\log_2n\rceil$.
Proof: if $k$ vertices are sources, all $n$ distinct out-neighborhoods
are subsets of the $n-k$ nonsources, so $n\le2^{n-k}$. Conversely,
take an extensional acyclic core on $m=\lceil\log_2n\rceil$ vertices
and choose $n-m$ unused subsets as new source neighborhoods. Since
$n>2^{m-1}$ for $m\ge1$, any resulting family covers the core, so it has
exactly $n-m$ sources. The case $n=1$ is immediate. The upper bound
already appears in Tomescu's 2011 thesis, printed page 29.

\paragraph{A182162.}
Let $T(n,k)$ denote the labeled count. For $q=\lceil\log_2n\rceil$,
\[
 T(n,n-q)=n!\,\mathrm{A001192}(q)\binom{2^q-q}{n-q}.
\]
More generally, for $d\ge0$ and $k=n-q-d\ge1$, formula
\eqref{sbd:eq:defect}, multiplied by $n!$, gives an exact $(d+1)$-term
expression. These are specializations of the marked-source sieve,
with the extremal case also admitting the direct Boolean-core bijection.

\paragraph{Boundary sequence for a possible separate entry.}
The unlabeled extremizer sequence begins
\[
 1,1,2,1,20,20,10,2,7128,8316,7128,4455,1980,594,108,9,\ldots.
\]
It satisfies $b_{2^m}^{(0)}=\mathrm{A001192}(m)$,
$\liminf(b_n^{(0)})^{1/n}=1$, and $\limsup(b_n^{(0)})^{1/n}=4$,
and it is not P-recursive. No assertion is made that this sequence
currently lacks an OEIS entry; a sequence-value and literature search
should precede any new-entry proposal.

\begin{writenote}[Added 3 October 2026, batch~85]
This appendix is draft text. Nothing in it was submitted to OEIS, by the
package's author (question~10 of Section~\ref{sbd:sec:questions}) or by the
intake, and this report submits nothing. Any proposal must keep
Tomescu's credit for the upper bound and is a human editor's decision after
independent review. The intake did not search OEIS for the boundary
sequence displayed above.
\end{writenote}

\section{Proof-status ledger}
\begin{longtable}{p{0.32\textwidth}p{0.61\textwidth}}
\toprule
Statement & Status and qualification\\
\midrule\endhead
Maximum-source upper bound & Known in Tomescu (2011); independently proved here.\\
Maximum formula and full support & Proved constructively here; supplies a proof of the displayed OEIS conjecture. No historical-first claim.\\
Finite collapse and rigidity & Classical; proofs included to fix semantics and labeling factors.\\
General source sieve & Existing enumeration; proof included and attribution retained.\\
Exact extremizer count & Proved by a classification and also by specialization of the sieve; no priority certification.\\
Finite-defect formula & Exact specialization with a proved $d+1$-term truncation.\\
Uniform covering bound & Proved by a corewise probability estimate.\\
Entropy law and all limit rates & Proved uniformly for each fixed defect using elementary bounds on $u_m$.\\
Non-P-recursiveness & Proved for every fixed defect, including labeled counts.\\
All-orders expansion & Proved away from $p=1$, with the core amplitude retained exactly.\\
Boundary binomial arc law & Proved in total variation with an explicit finite bound.\\
Numerical verification & Exact recurrence checks through 64 and exhaustive class enumeration through 7; not a substitute for the proofs.\\
Lean or Rocq verification & Not performed; no such certification claimed.\\
Natural boundary for $B_d$ & Not proved; proposed as a further question.\\
Historical priority of refinements & Not established by the available search.\\
\bottomrule
\end{longtable}

% [write] ragged right: the long URL and commit hash below left three underfull lines.
\begingroup\raggedright
\begin{thebibliography}{99}
\bibitem{oeis220}
OEIS Foundation, \emph{A182220: Largest number of sources of an extensional
acyclic digraph}, \url{https://oeis.org/A182220}.
Consulted October 3, 2026. The consulted entry labels the logarithmic
formula conjectural.

\bibitem{oeis162}
OEIS Foundation, \emph{A182162: Extensional acyclic digraphs by sources},
\url{https://oeis.org/A182162}. Includes Nathaniel Johnston's
marked-source-sieve Maple program, April 18, 2012.
Consulted October 3, 2026.

\bibitem{oeis1192}
OEIS Foundation, \emph{A001192: Number of full sets of size $n$},
\url{https://oeis.org/A001192}. Consulted October 3, 2026.

\bibitem{oeis161}
OEIS Foundation, \emph{A182161: Number of extensional acyclic digraphs on
$n$ labeled nodes}, \url{https://oeis.org/A182161}.
Consulted October 3, 2026.

\bibitem{peddicord}
Richard Peddicord, \emph{The number of full sets with $n$ elements},
Proceedings of the American Mathematical Society \textbf{13} (1962),
825--828. Bibliographic details checked through A001192 and Tomescu's
thesis; the original article was not independently inspected.

\bibitem{policriti}
Alberto Policriti and Alexandru I. Tomescu,
\emph{Counting extensional acyclic digraphs},
Information Processing Letters \textbf{111}(16) (2011), 787--791.
\url{https://doi.org/10.1016/j.ipl.2011.05.014}.
Publisher abstract inspected; closely related detailed arguments were
consulted in the following thesis.

\bibitem{tomescu}
Alexandru Ioan Tomescu, \emph{Sets as Graphs}, Ph.D. thesis,
Universit\`a degli Studi di Udine, December 2011.
\url{https://www.cs.helsinki.fi/u/tomescu/PhDThesis-AT.pdf}.
Relevant material: Chapter 2, Sections 2.1.1--2.1.2, printed pages 26--29;
source bound on page 29, rigidity on page 27, source deletion on page 28.
These sections were inspected in the available parsed text.

\bibitem{wagner}
Stephan Wagner, \emph{Asymptotic enumeration of extensional acyclic
digraphs}, Proceedings of ANALCO 2012, 1--8.
\url{https://doi.org/10.1137/1.9781611973020.1}.
Publisher abstract inspected. The author-hosted full-text link did not
respond during this research; no unchecked theorem from it is used in
our proofs.

\bibitem{repo}
Vladimir Reshetnikov, \emph{ProveIt}, repository snapshot
\url{https://github.com/VladimirReshetnikov/ProveIt}, commit
\code{6bf7f30d0352f7596e70928b3d4f304914075907},
inspected October 3, 2026. Used for project context and a limited
nonduplication search, not as a mathematical premise.
\end{thebibliography}
\endgroup
\end{document}
